%=======================================================================
%  IE 5311-001 (D01) Exploration Problem EP9
%  One page, single spacing, 10 pt, 1-inch margins, as the sheet requires.
%  Content past page 1 is not graded, so the page limit is hard.
%  Sections map onto the parts asked: what the cones are, the dual cone
%  and its characterisation, problems that use them, and conic duality.
%
%  IF IT SPILLS ONTO A SECOND PAGE, delete in this order until it fits.
%    1. the last sentence of section 5
%    2. the p = 7.3 row of the table in section 4
%    3. the sentence in section 3 beginning "Market impact"
%  Do not cut the AM-GM argument in section 2 or the table in section 4:
%  they are the dual cone and the duality the question asks about.
%=======================================================================
\documentclass[10pt]{article}

\usepackage[margin=1in]{geometry}
\usepackage{amsmath}
\usepackage{amssymb}
\usepackage{booktabs}
\usepackage{array}

\newcommand{\T}{^{\!\top}}
\newcommand{\Rp}{\mathbb{R}_+}

\pagestyle{empty}
\setlength{\parskip}{0.35em}
\setlength{\parindent}{0pt}

\begin{document}

{\large\bfseries EP9.\quad Power cones and their duals}\\[1pt]
{\small Lidivine Kengne \quad\textbar\quad IE 5311-001 (D01), Homework 6}
\vspace{0.3em}

\textbf{1. What they are.}
For $\alpha \in (0,1)$ the three-dimensional power cone is
\[
  \mathcal{P}_\alpha
  = \bigl\{ (x,y,z) \in \Rp \times \Rp \times \mathbb{R}
      \;:\; x^{\alpha} y^{1-\alpha} \ge |z| \bigr\}.
\]
It is a closed convex cone: the weighted geometric mean
$x^\alpha y^{1-\alpha}$ is concave on $\Rp^2$ and positively
homogeneous of degree one, so its hypograph is a cone and convex. The
general form replaces the two factors by $m$ of them and $z$ by a
vector, $\prod_j x_j^{\alpha_j} \ge \|z\|_2$ with
$\sum_j \alpha_j = 1$. Taking $\alpha = \tfrac12$ gives
$\sqrt{xy} \ge |z|$, which is exactly the rotated second-order cone ---
so power cones \emph{contain} the second-order cone as the single point
$\alpha = 1/2$ of a continuum.

\textbf{2. The dual cone.}
The dual is the same cone with its two nonnegative coordinates
rescaled:
\[
  \mathcal{P}_\alpha^{*}
  = \Bigl\{ (u,v,w) \in \Rp \times \Rp \times \mathbb{R}
      \;:\; \Bigl(\tfrac{u}{\alpha}\Bigr)^{\!\alpha}
            \Bigl(\tfrac{v}{1-\alpha}\Bigr)^{\!1-\alpha} \ge |w| \Bigr\}.
\]
Containment follows from weighted AM--GM. For nonnegative $x,y,u,v$,
\[
  \Bigl(\tfrac{u}{\alpha}\Bigr)^{\!\alpha}
  \Bigl(\tfrac{v}{1-\alpha}\Bigr)^{\!1-\alpha} x^{\alpha} y^{1-\alpha}
  = \Bigl(\tfrac{ux}{\alpha}\Bigr)^{\!\alpha}
    \Bigl(\tfrac{vy}{1-\alpha}\Bigr)^{\!1-\alpha}
  \;\le\; \alpha \cdot \tfrac{ux}{\alpha}
        + (1-\alpha)\cdot \tfrac{vy}{1-\alpha}
  = ux + vy ,
\]
so if $x^\alpha y^{1-\alpha} \ge |z|$ and the dual inequality holds,
then $ux + vy \ge |w||z| \ge -wz$, that is
$\langle (x,y,z),(u,v,w)\rangle \ge 0$. Sampling
(\texttt{ep9\_power.py}) checks both directions at
$\alpha = 0.2, 0.5, 0.75$: no pair drawn from
$\mathcal{P}_\alpha \times \mathcal{P}^*_\alpha$ had negative pairing,
and all $300$ sampled points outside the claimed dual admitted a cone
element they paired negatively with, so the characterisation is exact
and not merely sufficient. As a consistency check, an
$\alpha = \tfrac12$ program and the equivalent rotated second-order
cone program agree to $4.9 \times 10^{-8}$.

\textbf{3. Where they arise naturally.}
The headline case is the $p$-norm. For $p > 1$,
\[
  \|z\|_p \le t
  \quad\Longleftrightarrow\quad
  \exists\, s \ge 0: \;\; \sum_i s_i = t, \;\;
  |z_i| \le s_i^{1/p}\, t^{1-1/p} \;\; \text{for each } i,
\]
one power cone with $\alpha = 1/p$ per component. Geometric-mean
objectives $\prod_i x_i^{1/n} \ge t$ are the general form directly.
Market impact models in portfolio construction, where trading cost
grows like $|x|^{3/2}$, are a power cone with $\alpha = 2/3$.

\textbf{4. Conic duality on such a problem.}
Take $p$-norm regression, $\min_x \|Ax - b\|_p$, built from the cones
above. Its conic dual is
$\max_y\, b\T y$ subject to $A\T y = 0$ and $\|y\|_q \le 1$, where
$1/p + 1/q = 1$: the classical H\"older duality, recovered by the
conic machinery, with the rescaling by $1/\alpha$ and $1/(1-\alpha)$ in
$\mathcal{P}^*_\alpha$ being precisely what carries $p$ to its
conjugate exponent $q$. On a random instance with $60$ residuals and
$8$ unknowns:

\begin{center}\footnotesize
\setlength{\tabcolsep}{7pt}
\begin{tabular}{@{}rrrrrr@{}}
\toprule
$p$ & $q$ & power cone & SLSQP & H\"older dual & max gap\\
\midrule
1.50000 & 3.00000 & 3.57209424 & 3.57209420 & 3.57209420 & $3.7\times10^{-8}$\\
2.00000 & 2.00000 & 1.98539745 & 1.98539741 & 1.98539740 & $5.0\times10^{-8}$\\
3.00000 & 1.50000 & 1.13058808 & 1.13058807 & 1.13058802 & $6.2\times10^{-8}$\\
3.14159 & 1.46694 & 1.07640987 & 1.07640985 & 1.07640974 & $1.3\times10^{-7}$\\
7.30000 & 1.15873 & 0.62011110 & 0.62011109 & 0.62011107 & $3.1\times10^{-8}$\\
\bottomrule
\end{tabular}
\end{center}

Primal, dual and a general-purpose nonlinear solver agree throughout,
so the duality is not merely formal.

\textbf{5. Why not just use second-order cones.}
Because $\|\cdot\|_p$ is second-order cone representable only for
\emph{rational} $p$, and then by a tower whose size grows with the
continued-fraction expansion of $p$. Asked to build the $q$-norm of row
four by second-order cones, CVXPY reports that it approximates
$q = 1.46694$ with $11$ second-order cone constraints at an error of
$5.8 \times 10^{-8}$; row five costs $8$ constraints. One power cone
per residual is exact, fixed in size, and works for $p = \pi$ --- where
no tower exists at all.

\end{document}
